En Lenses, el F1 macro final es 1.000 para C4.5 y 1.000 para Random Forest. En Iris, el F1 macro final es 0.933 para C4.5 y 0.933 para Random Forest. En Cleveland, el F1 macro final es 0.719 para C4.5 y 0.836 para Random Forest. Estas cifras corresponden a la selección interna por validación; no al máximo de los barridos de prueba.

En Lenses, al pasar de cero a treinta por ciento de ausencias adicionales, C4.5 pasa de 1.000 a 0.230; Random Forest pasa de 0.619 a 0.352 de F1 macro medio. En Iris, al pasar de cero a treinta por ciento de ausencias adicionales, C4.5 pasa de 0.933 a 0.803; Random Forest pasa de 0.900 a 0.878 de F1 macro medio. En Cleveland, al pasar de cero a treinta por ciento de ausencias adicionales, C4.5 pasa de 0.703 a 0.790; Random Forest pasa de 0.885 a 0.830 de F1 macro medio. La comparación conserva parámetros y casos; no se supone que el efecto sea monótono ni que se repita en otros problemas.

En el barrido descriptivo de Cleveland, el árbol sin filtro obtiene 0.754 de tasa de acierto media, con desvío 0.043; con filtro obtiene 0.738, con desvío 0.043. Se informa la variante sin activarla por defecto. Ambas condiciones conservan los demás parámetros y las particiones. La comparación describe esta muestra y no selecciona el uso del filtro mediante los resultados de prueba.
